\documentclass[10pt,a4paper]{amsart}
\usepackage[margin=19mm]{geometry}
\usepackage{amsmath,amssymb,mathtools}
\usepackage[scr=rsfs,cal=euler]{mathalfa}
\usepackage{xfrac}
\usepackage[unicode,hidelinks,bookmarks=false]{hyperref}
\newcommand{\ie}{i.e.\ }
\newcommand{\cf}{cf.\ }
\newcommand{\resp}{respectively\ }
\newcommand{\Subsection}{\subsection}
\providecommand{\nobreakdash}{\nobreak\hbox{-}}
\setlength{\emergencystretch}{2em}
\multlinegap0pt
\usepackage{tikz}
\usepackage{amssymb}
\usepackage{stackengine}
\usepackage{enumerate}
\usepackage{xfrac}
\newtheorem{proposition}{Proposition}
\newtheorem{theorem}{Theorem}
\title{On quandle representations}
\author{Mohamad Maassarani}
\date{Source: arXiv:2605.12692v1 (CC BY 4.0)}
\begin{document}
\maketitle

\noindent\emph{Source and licence.} On quandle representations. Version arXiv:2605.12692v1 (CC BY 4.0), distributed by arXiv under the Creative Commons Attribution 4.0 International licence, http://creativecommons.org/licenses/by/4.0/, as recorded in the arXiv licence metadata for this exact version and shown on the abstract page https://arxiv.org/abs/2605.12692v1. Full licence notice: \texttt{licenses/arxiv-2605.12692-CC-BY-4.0.txt}.\par\smallskip

\noindent\textit{Editorial scope.} Counted unit: the article's theorem theorem (source0.tex lines 214--216). The article's complete original proof of it is delivered with it (source0.tex lines 217--234, one \texttt{proof} environment of 18 lines). No mathematics is written, repaired or re-derived: the statement and the proof are the article's own lines, in the article's own notation, and no retained line is substituted or re-flowed.  Retained uncounted as the source-local context the proof needs: lines 158--161.  Nothing was omitted from any retained span, so the package carries no omission record.  No retained line contains a citation, so the wrapper carries no bibliography and no bibliography entry.  No retained line contains a figure environment, an image-inclusion command or drawing code, so no image is delivered and the package declares no asset dependency.  Editorial formatting only, in addition: an AMS \texttt{amsart} preamble that restates the article's own declarations and macros used by the retained lines, namely the environments \texttt{proposition}, \texttt{theorem} on separate counters, together with the standard packages those lines need.  The retained environments print in this document as Proposition 1, Theorem 1 in order of appearance, whereas the article prints the same items with the numbering of its own full document.  The complete native source of the article as distributed in the arXiv e-print for this exact version is bundled unchanged as \texttt{sources/200437/source0.tex}.
\begin{proposition}\label{Z0} For $Q$ a finite quandle let $Z_0$ be the subgroup of $G(Q)$ generated by the elements $x^n$ for $x\in Q$ and with $n=\vert Inn(Q) \vert $.
\item[1)] The group $Z_0$ lies in the center of $G(Q)$
\item[2)] $Z_0$ has finite index in $G(Q)$.
\end{proposition}

\begin{theorem}
Let $Q$ be a finite quandle and $\rho : Q \to GL(V)$ be an irreducible finite dimensional complex representation of $Q$. $\rho$ is unitary for some inner product if and only if the determinant of every element of $\rho(Q)$ has modulus $1$.
\end{theorem}

\begin{proof}
If $\rho$ is unitary the elements of $\rho(Q)$ are isometries and hence have determinant with modulus $1$. We will prove the converse. Assume that the determinant of any element of $\rho(Q)$ has modulus $1$. Let $\rho_{G(Q)} : G(Q) \to GL(V)$ be the corresponding representation of $G(Q)$. Since the image of $\rho_{G(Q)}$ consist of products of elements of the image of $\rho$, the determinant of any element in the image of $\rho_{G(Q)}$ has modulus $1$. By proposition \ref{Z0}, we have a central extension $1 \to Z_0 \to G(Q) \to G(Q)/Z_0\to 1$ with $G(Q)/Z_0$ finite. Set $H=G(Q)/Z_0$. Let $s:H\to G(Q)$ be a section of $G(Q)\to H$ and $\langle-,-\rangle$ be an inner product of $V$. We define the inner product $\langle -,-\rangle_s$ on $V$ by :
$$ \langle v,w\rangle_s= \sum_{h\in H} \langle \rho_{G(Q)}(s(h))(v), \rho_{G(Q)}(s(h))(w)\rangle,$$
for $v,w\in V$. Take $g \in G$; $g$ can be written as $zs(h_1)$ for some $z\in Z_0$ and $h_1\in H$. We have that :
$$ \langle \rho_{G(Q)}(g)(v),\rho_{G(Q)}(g)(w)\rangle_s= \sum_{h\in H} \langle \rho_{G(Q)}(z_hs(hh_1))(v), \rho_{G(Q)}(z_hs(hh_1))(w)\rangle$$
where 
$$ z_h=zs(h)s(h_1)s(hh_1)^{-1},$$
lies in $Z_0$. But $Z_0$ lies in the center and the representation $\rho_{G(Q)}$ is irreducible with all elements in the image having a determinant of modulus $1$. Hence, $\rho_{G(Q)}(z_h)=\lambda_h id$ with $\lambda_h$ of modulus $1$. This proves that  :
\begin{equation*}
 \begin{split}
\langle \rho_{G(Q)}(g)(v),\rho_{G(Q)}(g)(w)\rangle_s  & = \sum_{h\in H} \langle \lambda_h\rho_{G(Q)}(s(hh_1))(v), \lambda_h\rho_{G(Q)}(s(hh_1))(w)\rangle\\
&=\sum_{h\in H}\lambda_h\overline{\lambda_h} \langle \rho_{G(Q)}(s(hh_1))(v), \lambda_h\rho_{G(Q)}(s(hh_1))(w)\rangle\\
&=\sum_{h\in H} \langle \rho_{G(Q)}(s(hh_1))(v), \rho_{G(Q)}(s(hh_1))(w)\rangle\\
&= \langle v,w\rangle_s
\end{split}
\end{equation*}
Therefore, the group representation $\rho_{G(Q)}$ is unitary with respect to $\langle -,-\rangle_s$ and so is $\rho$. We have proved the theorem.
\end{proof}

\end{document}
